\section*{Limitations}
\label{sec:limitations}

\paragraph{Scale and backbones.}
All experiments run on one laptop with compute tier S: a frozen \texttt{bge-small-en-v1.5} encoder, a frozen \texttt{flan-t5-small} answer head and a policy of about 0.3M parameters, trained on 3{,}000 questions with a 48{,}000-trajectory budget.
Absolute accuracy is therefore far below published HybridQA systems, which use much larger models, and our numbers are not comparable to them.
Nothing here shows how hop-level credit behaves for large language models, and we use a single backbone family.

\paragraph{Noisy gold paths.}
HybridQA has no annotated hop order. Our gold paths are heuristics built from a weak answer-node label, none is verified as clean, and the hop reward and path metrics inherit this noise.
A high strict validity means agreement with the heuristic, not necessarily correctness of reasoning.

\paragraph{Evidence from few runs and tuning.}
We use three seeds, a three-trial tuning grid per method (fewer than the six originally planned) and tune on seed 0 only, so seed variance is only partly captured and hyperparameter differences between methods may remain.
The reward weights are fixed and not tuned, and the row-evidence and feature changes were designed on development data after a first run performed poorly; all discarded runs are excluded and described in Appendix~\ref{app:changes}.

\paragraph{Benchmarks and OOD.}
We evaluate on HybridQA only. The OOD split holds out three table clusters within the same dataset and is not a cross-dataset test.
MultiHiertt, TAT-QA and FinQA were not run (Appendix~\ref{app:cut}).
Candidate generation caps accuracy: the gold first entity is not always among the $K=32$ hop-0 candidates.

\paragraph{Claims.}
Hypotheses whose status in the claims ledger is not ``supported'' are reported as such.
Unsupported or inconclusive hypotheses: H1 is not supported (\hop\ is not more accurate than answer-only GRPO, and seeds disagree in sign); H1c is not supported (the terminal-only variant is not less accurate, and seed variance is large); H2 is only partially supported because strict validity is low; H3 is inconclusive because the smaller relative OOD gap was not tested for significance and absolute OOD accuracy is low. With three seeds, seed-to-seed variance of the RL methods is large (up to about 6 EM points), which limits every RL comparison.
